\begin{titlepage}
\centering
{\Large Notas didácticas de revisión técnica\par}
\vspace{1.0cm}
{\huge\bfseries Un recorrido abierto por los artículos de investigación: cuatro ejes de la historia de la IA}\par
\vspace{0.5cm}
{\Large Zhang Xiaojun conversa con Xie Qingchi (谢青池)\par}
\vspace{0.35cm}
{\large Elaborado a partir del video público de Zhang Xiaojun Podcast, la transcripción generada en GPU, la hoja de ruta de artículos que figura en la descripción del video y diagramas conceptuales propios\par}
\vspace{0.3cm}
{\large 2025-10-28\par}
\vspace{0.85cm}
\includegraphics[width=0.88\textwidth,height=0.42\textheight,keepaspectratio]{cover.jpg}\par
\vfill
\begin{tcolorbox}[width=0.92\textwidth, colback=black!2!white, colframe=black!55, sharp corners]
\textbf{Título del video}: 117. Un recorrido abierto por los artículos de investigación: la historia completa de los cambios en paradigmas de modelos, Infra y datos, lenguaje y multimodalidad\par
\textbf{Canal del video}: Zhang Xiaojun Podcast\par
\textbf{Duración del video}: 04:22:37\par
\textbf{Enlace del video}: \href{\videourl}{\nolinkurl{\videourl}}\par
\textbf{Nota sobre los materiales}: YouTube no tiene subtítulos disponibles; el texto se elaboró a partir de la transcripción con faster-whisper large-v3 en CUDA, la hoja de ruta de 36 artículos incluida en la descripción de YouTube, la estructura de capítulos, la portada y diagramas didácticos propios. La versión de YouTube muestra una portada estática o una imagen fija, por lo que no se insertan fotogramas repetidos de los participantes; al obtener localmente los metadatos de la versión de Bilibili con pantalla compartida que se menciona en la descripción, el servidor devolvió HTTP 412, y la extracción de la presentación de Feishu no obtuvo el contenido, de modo que esta versión se basa en las fuentes locales confiables.\par
\textbf{Página del proyecto}: \href{\repourl}{\nolinkurl{\repourl}}\par
\end{tcolorbox}
\end{titlepage}

\tableofcontents
\newpage
